% =====================================================================
% Figure contract
% 核心结论：§3.4 用六条并列技术轴组织 CUA 综述，并把平台视为 Environment × Interface 的联合实例化。
% 图原型：原创中等档纯结构图；2×3 等尺寸卡片网格 + 底部整宽注释条。
% 证据层级：hero = 六色卡片网格；支撑 = 章节灰字与两行注释。
% 能不能更少？：不能；六张卡和两行注释均由用户逐项指定。
%
% Step ① 设计文档（Form A）
% 0 参考：沿用 cua-survey-fig1-overview.tex 的低饱和六色 palette、圆角与深灰文字。
% 1 领域：Computer-Use Agents 中文综述；英文轴名和技术术语保持原文。
% 2 格式：TikZ standalone；XeLaTeX + fontspec + xeCJK。
% 3 画布：W_ink 约 19.9 cm（含 standalone border 后 < 22 cm）；目标印刷宽度 19.9 cm；
%   目标印刷最小字号 7.5 pt，设计最小 CJK 字号 8.7 pt。
% 4 模块：六张 6.40×3.20 cm 卡片；标题 1 个；19.90×1.70 cm 注释条 1 个。
% 5 连线：无连线；卡头与卡身仅用同色细分隔线建立层级。
% 6 空间：列间距、行间距均 0.35 cm；卡头、分隔线、项目基线均从卡片锚点构造。
% 7 强调：轴名粗体；章节灰字；六张卡使用六种浅底色与对应中灰彩色描边。
% 8 密度：中等档纯分类结构，不嵌入无来源的数据可视化。
% 9 ASCII：
%   标题
%   ┌──────────┐ ┌──────────┐ ┌──────────┐
%   │Interface │ │Environment│ │Data      │
%   └──────────┘ └──────────┘ └──────────┘
%   ┌──────────┐ ┌──────────┐ ┌──────────┐
%   │Architecture││Learning │ │Evaluation│
%   └──────────┘ └──────────┘ └──────────┘
%   ┌──────────────────────────────────────┐
%   │ 两行整宽注释                         │
%   └──────────────────────────────────────┘
% 10 可视化决策：输入仅含分类文本，无数值/矩阵/曲线，故不添加 chart、heatmap 或公式。
%
% Pre-flight：P1 设计文档已写；P2 卡片用 positioning + anchors 构造；
% P3--P5 N/A（无连线/嵌入图）；P6 固定尺寸覆盖全部文字；P7 已读 lessons。
% =====================================================================

\documentclass[border=8pt]{standalone}
\usepackage{fontspec}
\usepackage{xeCJK}
\usepackage{xcolor}
\usepackage{tikz}
\setCJKmainfont{Heiti SC}
\setCJKsansfont{Heiti SC}
\usetikzlibrary{positioning,calc}

% Low-saturation palette shared with cua-survey-fig1-overview.tex.
\definecolor{ink}{HTML}{27364A}
\definecolor{muted}{HTML}{748092}
\definecolor{blueLine}{HTML}{6D879E}
\definecolor{blueFill}{HTML}{E7EEF4}
\definecolor{tealLine}{HTML}{648D8A}
\definecolor{tealFill}{HTML}{E5F0EE}
\definecolor{goldLine}{HTML}{9B875B}
\definecolor{goldFill}{HTML}{F2EEE3}
\definecolor{purpleLine}{HTML}{817594}
\definecolor{purpleFill}{HTML}{EEEAF2}
\definecolor{greenLine}{HTML}{71896E}
\definecolor{greenFill}{HTML}{EAF0E8}
\definecolor{roseLine}{HTML}{9A777B}
\definecolor{roseFill}{HTML}{F3E9EA}
\definecolor{noteFill}{HTML}{F1F3F6}
\definecolor{noteLine}{HTML}{B2BAC4}

\tikzset{
  every node/.style={
    font=\sffamily\fontsize{9.2}{11.2}\selectfont,
    text=ink
  },
  axiscard/.style={
    rectangle,
    rounded corners=5pt,
    line width=0.75pt,
    minimum width=6.40cm,
    minimum height=3.20cm,
    inner sep=0pt
  },
  cardhead/.style={
    anchor=north west,
    text width=5.72cm,
    align=left,
    inner sep=0pt,
    font=\sffamily\fontsize{10.3}{12.2}\selectfont
  },
  axisitem/.style={
    anchor=west,
    align=left,
    inner sep=0pt,
    font=\sffamily\fontsize{9.2}{11.2}\selectfont
  },
  notebar/.style={
    rectangle,
    rounded corners=4pt,
    draw=noteLine,
    fill=noteFill,
    line width=0.6pt,
    minimum width=19.90cm,
    minimum height=1.70cm,
    text width=18.98cm,
    align=left,
    inner xsep=0.42cm,
    inner ysep=0.24cm,
    font=\sffamily\fontsize{8.7}{12.8}\selectfont
  }
}

\begin{document}
\sffamily
\begin{tikzpicture}

% Card frames. Positioning coordinates make all row/column alignments exact.
\node[axiscard,anchor=north west,fill=blueFill,draw=blueLine] (card1) at (0,0) {};
\coordinate[right=0.35cm of card1.north east] (card2NW);
\node[axiscard,anchor=north west,fill=tealFill,draw=tealLine] (card2) at (card2NW) {};
\coordinate[right=0.35cm of card2.north east] (card3NW);
\node[axiscard,anchor=north west,fill=goldFill,draw=goldLine] (card3) at (card3NW) {};

\coordinate[below=0.35cm of card1.south west] (card4NW);
\node[axiscard,anchor=north west,fill=purpleFill,draw=purpleLine] (card4) at (card4NW) {};
\coordinate[right=0.35cm of card4.north east] (card5NW);
\node[axiscard,anchor=north west,fill=greenFill,draw=greenLine] (card5) at (card5NW) {};
\coordinate[right=0.35cm of card5.north east] (card6NW);
\node[axiscard,anchor=north west,fill=roseFill,draw=roseLine] (card6) at (card6NW) {};

% Card 1: Interface.
\node[cardhead] at ([xshift=0.34cm,yshift=-0.30cm]card1.north west)
  {{\bfseries Interface}\hfill{\color{muted}\fontsize{8.8}{10.4}\selectfont §4}};
\draw[blueLine!58,line width=0.45pt]
  ([xshift=0.34cm,yshift=-0.96cm]card1.north west) --
  ([xshift=-0.34cm,yshift=-0.96cm]card1.north east);
\node[axisitem] at ([xshift=0.45cm,yshift=-1.35cm]card1.north west)
  {{\color{blueLine}\textbullet}\hspace{0.55em}观测表示};
\node[axisitem] at ([xshift=0.45cm,yshift=-1.94cm]card1.north west)
  {{\color{blueLine}\textbullet}\hspace{0.55em}动作编码};
\node[axisitem] at ([xshift=0.45cm,yshift=-2.53cm]card1.north west)
  {{\color{blueLine}\textbullet}\hspace{0.55em}跨平台统一};

% Card 2: Environment.
\node[cardhead] at ([xshift=0.34cm,yshift=-0.30cm]card2.north west)
  {{\bfseries Environment}\hfill{\color{muted}\fontsize{8.8}{10.4}\selectfont §4}};
\draw[tealLine!58,line width=0.45pt]
  ([xshift=0.34cm,yshift=-0.96cm]card2.north west) --
  ([xshift=-0.34cm,yshift=-0.96cm]card2.north east);
\node[axisitem] at ([xshift=0.45cm,yshift=-1.35cm]card2.north west)
  {{\color{tealLine}\textbullet}\hspace{0.55em}reset·并行·fork};
\node[axisitem] at ([xshift=0.45cm,yshift=-1.94cm]card2.north west)
  {{\color{tealLine}\textbullet}\hspace{0.55em}verify·隔离·复现};

% Card 3: Data.
\node[cardhead] at ([xshift=0.34cm,yshift=-0.30cm]card3.north west)
  {{\bfseries Data}\hfill{\color{muted}\fontsize{8.8}{10.4}\selectfont §5}};
\draw[goldLine!58,line width=0.45pt]
  ([xshift=0.34cm,yshift=-0.96cm]card3.north west) --
  ([xshift=-0.34cm,yshift=-0.96cm]card3.north east);
\node[axisitem] at ([xshift=0.45cm,yshift=-1.35cm]card3.north west)
  {{\color{goldLine}\textbullet}\hspace{0.55em}任务与初始状态};
\node[axisitem] at ([xshift=0.45cm,yshift=-1.94cm]card3.north west)
  {{\color{goldLine}\textbullet}\hspace{0.55em}轨迹与 validator};

% Card 4: Architecture.
\node[cardhead] at ([xshift=0.34cm,yshift=-0.30cm]card4.north west)
  {{\bfseries Architecture}\hfill{\color{muted}\fontsize{8.8}{10.4}\selectfont §6}};
\draw[purpleLine!58,line width=0.45pt]
  ([xshift=0.34cm,yshift=-0.96cm]card4.north west) --
  ([xshift=-0.34cm,yshift=-0.96cm]card4.north east);
\node[axisitem] at ([xshift=0.45cm,yshift=-1.35cm]card4.north west)
  {{\color{purpleLine}\textbullet}\hspace{0.55em}end-to-end};
\node[axisitem] at ([xshift=0.45cm,yshift=-1.94cm]card4.north west)
  {{\color{purpleLine}\textbullet}\hspace{0.55em}compositional};

% Card 5: Learning.
\node[cardhead] at ([xshift=0.34cm,yshift=-0.30cm]card5.north west)
  {{\bfseries Learning}\hfill{\color{muted}\fontsize{8.8}{10.4}\selectfont §7}};
\draw[greenLine!58,line width=0.45pt]
  ([xshift=0.34cm,yshift=-0.96cm]card5.north west) --
  ([xshift=-0.34cm,yshift=-0.96cm]card5.north east);
\node[axisitem] at ([xshift=0.45cm,yshift=-1.35cm]card5.north west)
  {{\color{greenLine}\textbullet}\hspace{0.55em}SFT 与 RL};
\node[axisitem] at ([xshift=0.45cm,yshift=-1.94cm]card5.north west)
  {{\color{greenLine}\textbullet}\hspace{0.55em}自我改进与搜索};

% Card 6: Evaluation.
\node[cardhead] at ([xshift=0.34cm,yshift=-0.30cm]card6.north west)
  {{\bfseries Evaluation}\hfill{\color{muted}\fontsize{8.8}{10.4}\selectfont §8}};
\draw[roseLine!58,line width=0.45pt]
  ([xshift=0.34cm,yshift=-0.96cm]card6.north west) --
  ([xshift=-0.34cm,yshift=-0.96cm]card6.north east);
\node[axisitem] at ([xshift=0.45cm,yshift=-1.35cm]card6.north west)
  {{\color{roseLine}\textbullet}\hspace{0.55em}程序化 verifier};
\node[axisitem] at ([xshift=0.45cm,yshift=-1.94cm]card6.north west)
  {{\color{roseLine}\textbullet}\hspace{0.55em}交互式 verifier};
\node[axisitem] at ([xshift=0.45cm,yshift=-2.53cm]card6.north west)
  {{\color{roseLine}\textbullet}\hspace{0.55em}LLM judge};

% Full-width note bar.
\coordinate[below=0.40cm of card4.south west] (noteNW);
\node[notebar,anchor=north west] (note) at (noteNW)
  {平台（Web / Mobile / Desktop / Hybrid）= Environment × Interface 两轴的联合实例化\\
   §9 产业与 §10 开放问题不是第七、八条轴，而是六轴结论向部署约束与研究议程的投射};

% Title is placed last so its baseline is independent of card geometry.
\node[
  anchor=south west,
  inner sep=0pt,
  font=\sffamily\fontsize{14.2}{17.0}\selectfont\bfseries,
  text=ink
] at ([yshift=0.58cm]card1.north west) {六条技术轴（§3.4）};

\end{tikzpicture}
\end{document}
